\par\noindent\textbf{(1)} 电子电荷为 \(-e\)、质量 \(m_e\)，轨道磁矩 \(\boldsymbol\mu_L=-e\boldsymbol L/(2m_e)\)。外场能量 \(H_Z=-\boldsymbol\mu_L\cdot\boldsymbol B=(eB/(2m_e))L_z=\mu_B B(L_z/\hbar)\)，其中Bohr磁子 \(\mu_B=e\hbar/(2m_e)\)。
\par\noindent\textbf{(2)} 因为 \(L_z|l,m_l\rangle=m_l\hbar|l,m_l\rangle\)，一阶能量修正 \(\Delta E_{m_l}=\mu_BBm_l\)，\(m_l=-l,-l+1,\ldots,l\)。原来 \(2l+1\) 重磁量子数简并裂为等间距能级，间距 \(\mu_BB\)。
\par\noindent\textbf{(3)} 电偶极算符是角动量为1的球张量，选择规则 \(\Delta l=\pm1\)、\(\Delta m_l=0,\pm1\)。发射或吸收光子能量的磁场修正 \(\Delta(\hbar\omega)=\mu_BB(m'_l-m_l)=0,\pm\mu_BB\)，形成正常Zeeman三重偏移，而不是任意 \(2l+1\) 条等间距可见谱线。
\par\noindent\textbf{(4)} 若存在非零自旋和旋轨耦合，弱磁场好量子数应为 \(J,m_J\)，磁移动 \(\Delta E=\mu_Bg_Jm_JB\)，其中 \(g_J=1+[J(J+1)+S(S+1)-L(L+1)]/[2J(J+1)]\)。当电子自旋参与时可出现反常Zeeman结构，不能无条件使用 \(g=1\)。
